% !TEX root = ../../../main.tex

\section{白盒单元测试}

\subsection{白盒单元测试设计}

白盒单元测试主要面向 HearBridge 后端项目展开。

后端是移动端、管理端和 Python 识别服务之间的业务协调中心，
其内部逻辑直接影响用户登录状态、资源管理、样本管理、训练任务、模型发布和识别结果返回等关键流程。

因此，
本文将白盒测试重点放在后端 Controller 层、Service 层和 Interceptor 层，
通过自动化测试检查接口入口、参数校验、业务状态判断、异常分支处理、认证拦截和外部依赖调用是否符合预期。

当前后端测试代码位于 \texttt{src/test/java/com/yezhen/hearbridge/backend} 目录下，
按照 \texttt{controller}、\texttt{service}、\texttt{interceptor} 等目录组织。

测试工程共包含 22 个 Java 测试类，
其中 Controller 测试类 10 个，
Service 测试类 10 个，
Interceptor 测试类 1 个，
Spring Boot 应用上下文测试类 1 个。

Controller 层测试主要用于验证接口入口层的请求参数接收、服务调用和响应封装；
Service 层测试主要用于验证核心业务逻辑、状态流转和异常处理；
Interceptor 层测试主要用于验证统一认证拦截、公开接口放行、App 端接口鉴权、管理端接口鉴权和 401 响应写出等安全相关逻辑。

测试实现主要基于 JUnit 5、Mockito 和 Spring Boot Test。

对于 Service 层和 Interceptor 层测试，
系统使用 Mockito 模拟 Mapper、Redis、MinIO、Python 服务客户端或请求响应对象，
避免测试结果受到真实数据库、缓存、文件存储和外部服务运行状态影响。

对于 Controller 层测试，
系统通过 MockMvc 模拟 HTTP 请求，
并检查接口调用后的状态码、响应字段和异常返回结构是否符合系统约定。

覆盖率统计采用 IntelliJ IDEA 自带 Coverage 分析工具完成，
不额外引入 JaCoCo 覆盖率插件。

从实际测试文件来看，
Controller 层测试覆盖了管理端认证、文件上传、手势分类、手势资源、手势样本、样本训练、手势搜索、模型版本、模型产物文件和句子视频语义修正等接口入口；
Service 层测试覆盖了管理端认证、移动端认证、手势分类、手势资源、样本管理、样本文件流转、模型版本管理、模型发布、句子视频语义修正和相关业务异常处理；
Interceptor 层测试覆盖了系统统一鉴权逻辑。

测试对象与当前 HearBridge 的后台管理闭环、移动端认证链路和句子视频识别语义修正链路保持一致，
能够较好支撑系统关键后端逻辑的正确性验证。

表~\ref{tab:chapter07-white-box-design} 给出了白盒单元测试的主要设计内容。

\begin{longtable}{L{0.20\textwidth} L{0.38\textwidth} L{0.34\textwidth}}
  \caption{白盒单元测试设计表}
  \label{tab:chapter07-white-box-design}\\
  \toprule
  测试层次 & 实际测试对象 & 测试重点 \\
  \midrule
  Controller 层 & 管理端认证、文件上传、手势分类、手势资源、手势样本、样本训练、模型版本、模型文件、手势搜索、语义修正等控制器测试 & 验证接口入口参数、服务调用、响应结构和异常返回是否符合接口约定 \\
  Service 层 & 管理端认证、移动端认证、手势分类、手势资源、样本管理、样本文件流转、模型版本管理、模型发布、语义修正等服务测试 & 验证业务状态判断、数据组装、异常分支、外部依赖调用和结果封装逻辑 \\
  Interceptor 层 & \texttt{AuthInterceptorTest} & 验证公开接口放行、App 接口鉴权、管理端接口鉴权、token 解析和 401 响应写出逻辑 \\
  应用上下文 & Spring Boot 应用上下文测试 & 验证后端应用基础上下文能够正常加载 \\
  \bottomrule
\end{longtable}

测试用例设计遵循正常流程与异常流程并重的原则。

正常流程用于验证管理员登录、用户登录、分类维护、资源维护、样本管理、模型发布、语义修正等主路径能够正确执行；
异常流程用于验证密码错误、资源不存在、参数缺失、状态不匹配、重复操作、token 缺失、token 无效、外部服务响应异常等情况下系统能否返回合理错误信息。

通过这种方式，
白盒测试不仅检查功能能否执行，
也检查异常情况下是否能够被正确拦截和反馈。

\subsection{白盒单元测试结果}

在完成白盒测试设计后，
本文对 HearBridge 后端项目执行自动化单元测试。

测试入口集中在 Maven 项目的 \texttt{src/test/java} 目录下，
测试执行与覆盖率统计主要通过 IntelliJ IDEA 自带的测试运行器和 Coverage 分析功能完成，
用于查看测试通过情况以及 Controller、Service、Interceptor 等后端代码包的覆盖情况。

图~\ref{fig:chapter07-white-box-test-success} 展示了后端白盒单元测试全部通过的运行结果。

\WhuAutoFigure{0.86\textwidth}{0.45\textheight}{figures/chapter07/white-box-test-success.png}{后端白盒单元测试全部通过截图}{fig:chapter07-white-box-test-success}

从测试执行结果来看，
当前后端自动化测试共执行 108 个测试用例，
全部测试均通过，
测试总耗时约 6.771 秒。

测试范围覆盖了管理端认证、文件上传、手势分类、手势资源、手势样本、样本训练、手势搜索、模型版本管理、模型产物发布、移动端用户认证、统一认证拦截以及句子视频语义修正等核心模块。

表~\ref{tab:chapter07-white-box-result} 给出了白盒单元测试执行结果统计。

\begin{longtable}{L{0.26\textwidth} R{0.18\textwidth} L{0.46\textwidth}}
  \caption{白盒单元测试执行结果统计表}
  \label{tab:chapter07-white-box-result}\\
  \toprule
  统计项 & 测试结果 & 说明 \\
  \midrule
  Java 测试类数量 & 22 个 & 包含 Controller、Service、Interceptor 和应用上下文测试 \\
  自动化测试用例数 & 108 个 & IDEA 测试运行器统计结果 \\
  测试通过数量 & 108 个 & 所有测试均执行通过 \\
  测试失败数量 & 0 个 & 未出现断言失败或运行异常 \\
  测试总耗时 & 约 6.771 秒 & 本地开发环境运行结果 \\
  主要测试范围 & Controller 层、Service 层、Interceptor 层 & 覆盖接口入口、核心业务逻辑与认证拦截逻辑 \\
  \bottomrule
\end{longtable}

图~\ref{fig:chapter07-white-box-coverage} 展示了 IntelliJ IDEA Coverage 对后端 Java 代码的覆盖率统计结果。

\WhuAutoFigure{0.88\textwidth}{0.48\textheight}{figures/chapter07/white-box-coverage.png}{后端白盒单元测试覆盖率截图}{fig:chapter07-white-box-coverage}

表~\ref{tab:chapter07-backend-coverage} 给出了后端代码覆盖率统计结果。

\begin{longtable}{L{0.24\textwidth} R{0.18\textwidth} R{0.18\textwidth} R{0.20\textwidth} R{0.18\textwidth}}
  \caption{后端代码覆盖率统计表}
  \label{tab:chapter07-backend-coverage}\\
  \toprule
  代码范围 & 类覆盖率 & 方法覆盖率 & 行覆盖率 & 分支覆盖率 \\
  \midrule
  项目总体 & 58\%（41/70） & 69\%（189/272） & 58\%（1058/1806） & 36\%（351/963） \\
  \texttt{controller} 包 & 100\%（13/13） & 75\%（43/57） & 67\%（73/108） & 25\%（3/12） \\
  \texttt{service} 包 & 88\%（16/18） & 69\%（116/167） & 60\%（835/1377） & 43\%（304/691） \\
  \texttt{interceptor} 包 & 100\%（2/2） & 100\%（10/10） & 92\%（98/106） & 66\%（32/48） \\
  \texttt{config} 包 & 75\%（6/8） & 76\%（10/13） & 48\%（36/74） & 26\%（10/38） \\
  \texttt{util} 包 & 0\%（0/2） & 0\%（0/14） & 0\%（0/124） & 0\%（0/170） \\
  \bottomrule
\end{longtable}

从覆盖率结果可以看出，
当前白盒测试已经覆盖后端的主要业务层和接口入口层。

其中，
Controller 包类覆盖率达到 100\%，
说明当前后端主要接口入口均已纳入自动化测试范围；
Service 包行覆盖率达到 60\%，
覆盖了手势资源管理、样本管理、训练管理、模型版本发布、移动端认证和语义修正等核心业务逻辑；
Interceptor 包行覆盖率达到 92\%，
能够较好支撑系统认证拦截与权限校验相关测试说明。

同时，
覆盖率结果也反映出当前测试仍存在一定改进空间。

项目总体分支覆盖率为 36\%，
主要原因是部分异常分支、配置分支、工具类分支和外部服务调用失败场景尚未完全覆盖；
\texttt{util} 包当前覆盖率为 0\%，
说明版本号处理、路径处理或通用工具方法等辅助逻辑尚未纳入单元测试。

由于该类代码不直接构成主要业务流程，
当前阶段对系统核心功能验证影响较小，
但后续可补充工具类边界测试和异常路径测试， 
以进一步提升整体覆盖率。

总体来看，
HearBridge 后端白盒测试已经覆盖系统核心业务链路，
尤其是后台管理、模型发布、移动端认证、统一认证拦截和语义修正等关键模块。

测试结果能够支撑毕业设计阶段对系统正确性和可靠性的验证要求。

后续若继续完善，
可优先补充工具类、配置类和外部服务异常分支相关测试。